\subsubsection{Diffuse product poles}
\label{subsec:diffuse-product-polar-obstruction}

Polar curvature on every modification still misses some ordinary models. A maximum of two diffuse curve potentials has a very small joint pole set. Its first polar currents vanish, while the bounded truncations retain positive top-degree mass near that set.

\begin{lemma}\label{lem:diffuse-product-truncation-mass}
Let $\Delta\subset\C$ be the unit disc. There are a compact set $K\Subset\Delta$ of Hausdorff dimension zero and a negative subharmonic function $g$ on $\Delta$, smooth and harmonic off $K$, with
\[
 \{g=-\infty\}=K,\qquad \nu(g,a)=0\quad(a\in\Delta),
 \qquad g|_{\partial\Delta}=0.
\]
Put $D=\Delta^2$, $S=K\times K$ and
\[
 h(z,w)=\max\{g(z),g(w)\},\qquad
 h_L=\max\{h,-L\}\quad(L>0).
\]
The function $h$ is locally Lipschitz off $S$ and its bounded local Monge--Amp\`ere measure there is zero. The measures $(\ddc h_L)^2$ have one common finite mass $M>0$. Their supports are compact and approach $S$ as $L\to\infty$.
\end{lemma}
\begin{proof}
Use the binary Cantor set and nonatomic probability $\sigma$ of \cref{ex:diffuse-polar-zero-model}, placed in a small real interval in $\Delta$. Its level-$k$ intervals have length $r_k=\exp(-4^{k+1})$ and mass $2^{-k}$. For every $s>0$,
\[
 2^kr_k^s\longrightarrow0,\qquad
 4^k(\sqrt2r_k)^s\longrightarrow0.
\]
These covers show that both $K$ and $S$ have Hausdorff dimension zero. Define the disc Green potential
\[
 g(z)=\int_\Delta
       \log\left|\frac{z-y}{1-\bar yz}\right|^2\,\mathrm d\sigma(y).
\]
Then $\ddc g=\sigma$, with the normalization of this book. The Green kernel is nonpositive, zero on the boundary, and smooth away from its pole. On a fixed smaller disc it is at most $\log|z-y|^2+C$. Integration over the level-$k$ interval containing $z\in K$ gives
\[
 g(z)\leq2^{-k}(2\log r_k+C)\longrightarrow-\infty.
\]
Outside $K$ the potential is finite, smooth and harmonic. Its Lelong number at $a$ is $\sigma(\{a\})=0$. Since $K$ stays away from the boundary, $g$ is smooth near the boundary and tends uniformly to zero there.

At a point of $D\setminus S$, either both curve potentials are harmonic locally, or the finite one dominates the other near its pole. Thus $h$ is locally a maximum of at most two smooth pluriharmonic functions and is locally Lipschitz. Such a maximum has zero Monge--Amp\`ere measure in dimension two: the decreasing smooth approximations
\[
 \varepsilon\log(e^{a/\varepsilon}+e^{b/\varepsilon})
\]
have complex Hessian of rank at most one when $a,b$ are pluriharmonic. Bounded decreasing convergence in \cref{thm:bt-input} gives the assertion for their maximum.

The bounded truncation $h_L$ is continuous: near a pole of either curve potential its branch is suppressed by $-L$. The same argument for the three branches $g(z),g(w),-L$ shows that $(\ddc h_L)^2$ is carried by
\[
 Q_L=\{(z,w):g(z)=g(w)=-L\}.
\]
At a point where one branch is strictly smaller, at most two pluriharmonic branches remain. Near a pole that branch is suppressed by $-L$. The set $Q_L$ is compact in $D$, so \eqref{eq:local-cln-range} makes the mass finite.

For $L,L'>0$, the functions $h_L,h_{L'}$ agree near $\partial D$, since $h$ tends uniformly to zero there. Bounded integration by parts gives
\[
 (\ddc h_L)^2-(\ddc h_{L'})^2
       =\ddc\bigl((h_L-h_{L'})(\ddc h_L+\ddc h_{L'})\bigr).
\]
The current on the right has compact support. Its integral is zero, so all masses equal a number $M$. To see $M>0$, fix $L$ and choose a sufficiently small $\varepsilon>0$ such that
\[
 U=\{h_L<\varepsilon(|z|^2+|w|^2-3)\}
\]
is nonempty. It is relatively compact in $D$, because $h_L=0$ on the boundary and the quadratic is strictly negative there. The bounded comparison \eqref{eq:interior-power-comparison} gives
\[
 0<\int_U\bigl(\ddc\varepsilon(|z|^2+|w|^2)\bigr)^2
       \leq\int_U(\ddc h_L)^2\leq M.
\]
Finally, $g$ is bounded on compact subsets of $\overline\Delta\setminus K$. Its deep level sets therefore approach $K$, and $Q_L$ approaches $S$.
\end{proof}

\Needspace{12\baselineskip}
\begin{proposition}\label{prop:all-modifications-diffuse-polar-blindness}\index{Polar curvature!diffuse product poles}\index{Model potential!all modifications and separation}
On $X=\mathbb P^1\times\mathbb P^1$ with a K\"ahler form $\omega$, there are a normalized ordinary model $p\ne0$ and $\delta>0$ such that $p,0\in\mathcal M_\delta(\omega)$, all multiplier ideals of $p$ are trivial, and
\[
 \mathbf1_{\{p\circ\pi=-\infty\}}
       \bigl(\pi^*\omega+\ddc(p\circ\pi)\bigr)=0
\]
for every smooth modification $\pi:Y\to X$. Consequently, for every allowed family and every choice of Gauduchon forms,
\[
 \delta_\omega^{\mathrm{pol}}(p,0)=0,\qquad d_\omega(p,0)>0.
\]
The polar comparison fails to separate ordinary models of positive lower volume. Its analytic closure contains the nonideal model $p$.
\end{proposition}
\begin{proof}
Place the bidisc of the lemma in an affine chart of $X$. Choose a smooth cutoff $0\leq\chi\leq1$, compactly supported in that bidisc and equal to one on a neighborhood of $S$. Extend $H=\chi h$ by zero. The function is genuine and quasi-plurisubharmonic: outside $S$, the gradients of $h$ are locally bounded, so the terms containing derivatives of $\chi$ in
\[
 \ddc(\chi h)=\chi\ddc h
       +h\ddc\chi+\mathrm d\chi\wedge\dc h
       +\mathrm dh\wedge\dc\chi
\]
are bounded coefficient functions on their compact support. The first term is positive. Hence $\ddc H\geq-C\omega$ for some $C>0$. Take $0<t<1/C$ and put $u=tH$. Then $u\leq0$, $\sup_Xu=0$, its pole set is $S$, and
\[
 \omega+\ddc u\geq(1-tC)\omega.
\]

Set $u_k=\max(u,-k)$. The bounded currents $\omega+\ddc u_k$ have total Monge--Amp\`ere mass $A=\int_X\omega^2$ by closed-background Stokes. Near $S$, where $\chi=1$,
\[
 (\omega+\ddc u_k)^2
       =(\omega+t\ddc h_{k/t})^2
       \geq t^2(\ddc h_{k/t})^2.
\]
For large $k$, the entire support of the measure on the right lies in this neighborhood. Those measures have mass $t^2M$ and approach $S$. Take a common weakly convergent subsequence of these local measures and the global bounded-product measures. The global limit has mass at least $t^2M$ on $S$. Off $S$, the function $u$ is locally bounded, and $u_k=u$ eventually on every compact set. Thus the global limit there is precisely the non-pluripolar measure of $u$. It follows that
\[
 (1-tC)^2A
       \leq\int_X\mathrm{MA}_\omega(u)
       \leq A-t^2M<A.
\]
The lower bound follows from the strict current inequality and bounded local products off the Lebesgue-null set $S$.

In a closed K\"ahler background, ordinary bounded equivalence preserves non-pluripolar mass. Therefore both singularity volumes of $u$ equal this positive number. By \cref{thm:ordinary-envelope-volumes}, $p=P_\omega[u]$ is a normalized ordinary model with the same volumes. In particular $p\ne0$, and $d_\omega(p,0)=A-\overline{\operatorname{vol}}_\omega(p)>0$. Taking half its positive lower volume for $\delta$ puts both models in $\mathcal M_\delta(\omega)$. Also $u\leq p\leq0$, so the pole set of $p$ is contained in $S$.

Fix a modification $\pi$ of these smooth compact surfaces. Its exceptional locus is a finite union of curves $E_i$, and it is an isomorphism away from them and their finitely many image points. The polar current $R$ of $p\circ\pi$ is closed and positive by \cref{prop:polar-curvature-closedness}. Outside the exceptional curves it is carried by the inverse image of $S$, which has Hausdorff dimension zero in each relatively compact isomorphism chart. By \cref{cor:closed-current-hausdorff-null}, $R$ vanishes there. The analytic support theorem now gives
\[
 R=\sum_i b_i[E_i],\qquad b_i\geq0.
\]
See \cite[Chapter III, Corollaries 2.11 and 2.14, pp.~141--143]{DemCADG}. Since $R\leq\pi^*\omega+\ddc(p\circ\pi)$, its coefficients satisfy $b_i\leq\nu(p\circ\pi,E_i)$.

If $\pi(E_i)\notin S$, the lower bound by $u\circ\pi$ makes $p\circ\pi$ bounded near $E_i$, so $b_i=0$. If $\pi(E_i)=(a_1,a_2)\in S$, a neighborhood of that point has $u=t\max(g(z),g(w))$. At a general point of $E_i$, the first coordinate satisfies
\[
 z\circ\pi-a_1=\zeta^m f,\qquad m\geq1,\quad f\ne0.
\]
A local holomorphic root of the unit $f$ makes this expression $\zeta'^m$. Circular means give
\[
 \nu(g(a_1+\zeta'^m),E_i)=m\nu(g,a_1)=0.
\]
Since $p\circ\pi\geq u\circ\pi\geq t\,g(z\circ\pi)$ there, again $b_i=0$. Thus $R=0$ for every modification.

The same local comparison gives zero Lelong numbers for $u$, and hence for $p\geq u$, at every point. Skoda's local integrability theorem, recalled in \cite{SGPT}, makes $\mathrm e^{-sp}$ locally integrable for every $s>0$. Thus all multiplier ideals of $p$ are trivial. Its ideal envelope is therefore the zero model. Since $p\ne0$, it is nonideal. The zero model is analytic, and $\delta_\omega^{\mathrm{pol}}(p,0)=0$ puts $p$ in its polar-comparison closure. This proves the final assertions.
\end{proof}
